\documentclass[12pt]{article}
\usepackage[utf8]{inputenc}
\usepackage[T1]{fontenc}
\usepackage{hyperref}
\title{Termanu Language Resources for K-12 Curriculum Development}
\author{Compiled by Claude Sonnet 5 (claude-sonnet-5) for Andr\'as Kornai}
\date{}
\begin{document}
\maketitle
\section{Introduction}
Termanu (ISO 639-3 twu) is a Central Malayo-Polynesian language of Rote island, Indonesia, with about 30,000 speakers. Searches found few modern teaching materials.
\section{Primary Educational Resources}
None found.
\section{Linguistic and Reference Materials}
\begin{itemize}
\item Wikipedia on Termanu, which cites Jonker's 1908 Rottineesch-Hollandsch Woordenboek: \url{https://en.wikipedia.org/wiki/Termanu_language}
\item An OAPEN open-access book returned by the search (its title was not shown, so its relevance is unverified): \url{https://library.oapen.org/bitstream/20.500.12657/52033/1/book.pdf}
\item PARADISEC collection TWU2018: \url{https://catalog.paradisec.org.au/repository/TWU2018}
\end{itemize}
\section{Religious and Community Resources}
Scripture Earth lists the Gospel of Mark in Termanu: \url{https://www.scriptureearth.org/00eng.php?iso=twu}
\section{Digital and Audio Resources}
\begin{itemize}
\item Global Recordings Network, Termanu: \url{https://globalrecordings.net/en/language/twu}
\item Global Recordings Network, related page: \url{https://globalrecordings.net/en/language/25547}
\end{itemize}
\section{Curriculum Development Context}
Termanu is vital in the community but absent from schools. A curriculum would have to be built from the PARADISEC recordings, the historical dictionary and new documentation work.
\end{document}